\section{Methods}
\label{sec:methods}
%% ========================================================================

\subsection{Development Process and Timeline}

The repository and this paper were constructed using an AI-assisted development methodology. Kevin Kawchak served as the prompt creator, system architect, and provided edits to the generated paper, directing AI code generation through structured prompts documented in \texttt{prompts.md}.

The primary AI system used for repository generation and this paper was \textbf{Anthropic Claude Code Opus 4.6}. Peer review was conducted using \textbf{OpenAI ChatGPT 5.4 Thinking}, with review documents archived under the \texttt{peer-review/} directory.

\subsubsection{Repository Evolution}

The system was developed across four repositories in sequence. Table~\ref{tab:repo-timeline} summarizes the development timeline:

\begin{table}[H]
\centering
\caption{Repository development timeline and evolution.}
\label{tab:repo-timeline}
\scriptsize
\begin{tabularx}{\textwidth}{@{}L{4.7cm}L{1.6cm}R{0.7cm}L{5.2cm}@{}}
\toprule
\textbf{Repository} & \textbf{Date Range} & \textbf{Days} & \textbf{Focus} \\
\midrule
\texttt{physical-ai-oncology-trials}~\cite{kawchak2026physicalai} & Feb 2026 & 14 & USL scoring, patient instructions, single-site Physical AI framework \\
\texttt{pai-oncology-trial-fl}~\cite{kawchak2026paifl} & Feb--Mar 2026 & 10 & Federated learning, privacy modules, regulatory compliance \\
\texttt{mcp-pai-oncology-trials}~\cite{kawchak2026trialmcp} & Mar 2026 & 7 & MCP server proof of concept, 5-server architecture, 23 tools \\
\texttt{national-mcp-pai-oncology-trials}~\cite{kawchak2026national} & Mar 2026 & 5 & National standard, SDKs, conformance suite, governance \\
\bottomrule
\end{tabularx}
\end{table}

%% NOTE: This is chunk 02 of 03 - see README.md for reconstruction instructions
%% Content continues from 01_preamble_introduction.tex and continues in 03_discussion_conclusion.tex